% Complete dependency blueprint for arXiv:2505.14338v3.
% Paper-specific results are intentionally not marked \leanok/\mathlibok.
% External results whose proofs are cited but not reproduced are marked
% \notready rather than being silently treated as available.

\newcommand{\R}{\mathbb{R}}
\newcommand{\CPWL}{\operatorname{CPWL}}
\newcommand{\ReLU}{\operatorname{ReLU}}
\newcommand{\MAX}{\operatorname{MAX}}
\newcommand{\PolyC}{\mathcal{P}}
\newcommand{\Tspace}{\mathcal{T}}

\chapter{Neural-network depth and the main results}

\begin{definition}[Rectified linear unit]
  \label{def:relu-scalar}
  The rectified linear unit is the scalar function
  $\ReLU(t)=\max\{0,t\}$.  For a vector, $\ReLU$ denotes coordinatewise
  application of this function.
\end{definition}

\begin{definition}[ReLU network and hidden-layer depth]
  \label{def:relu-network}
  \uses{def:relu-scalar}
  Let $n_0,\ldots,n_{k+1}$ be positive integers and let
  $T^{(i)}:\R^{n_{i-1}}\to\R^{n_i}$ be affine for
  $i=1,\ldots,k+1$.  The associated ReLU network computes
  \[
    T^{(k+1)}\circ\ReLU\circ T^{(k)}\circ\cdots\circ
    \ReLU\circ T^{(1)}.
  \]
  It has $k$ hidden layers.  Its size is the total number of hidden neurons,
  equivalently the sum $n_1+\cdots+n_k$ in the layered presentation.
\end{definition}

\begin{definition}[Continuous piecewise-linear function]
  \label{def:cpwl-function}
  A function $f:\R^n\to\R$ is continuous piecewise linear (CPWL) if it is
  continuous and there are finitely many polyhedra covering $\R^n$ such that
  the restriction of $f$ to each polyhedron is affine.
\end{definition}

\begin{lemma}[ReLU networks compute CPWL functions]
  \label{lem:relu-network-is-cpwl}
  \uses{def:relu-network, def:cpwl-function}
  Every scalar-valued function computed by a finite ReLU network is CPWL.
\end{lemma}
\begin{proof}
  \uses{def:relu-scalar, def:relu-network, def:cpwl-function}
  Each affine coordinate function is CPWL.  If $g$ is CPWL, then the finitely
  many cells on which $g$ is affine can be cut further by the affine
  inequalities $g\geq0$ and $g\leq0$; on the resulting finite polyhedral
  refinement, $\ReLU\circ g$ equals either $g$ or zero and is therefore CPWL.
  A finite tuple of CPWL functions admits the common refinement of their
  finitely many cell decompositions, and postcomposition of that tuple with an
  affine map is affine on every cell of the refinement.  Starting with
  $T^{(1)}$ and applying these two observations alternately through all hidden
  layers proves the claim by induction on $k$.
\end{proof}

\section{The ternary depth recursion}

\begin{definition}[The functions $T_{a,b}$]
  \label{def:t-function}
  \uses{def:max-function}
  For integers $a,b\geq0$, define $T_{a,b}:\R^{4a+b}\to\R$ by
  \[
  T_{a,b}(x)=\max\left(
    \left\{\max(x_{4j+1},x_{4j+2})+
             \max(x_{4j+3},x_{4j+4}):0\leq j<a\right\}
    \cup\{x_{4a+1},\ldots,x_{4a+b}\}\right).
  \]
  Thus $T_{0,b}=\MAX_b$.
\end{definition}

\begin{definition}[The vector spaces $\Tspace_{a,b}(N)$]
  \label{def:t-space}
  \uses{def:t-function}
  For an ambient input dimension $N$, let $\Tspace_{a,b}(N)$ be the real
  vector space spanned by all functions $T_{a,b}\circ L$, where
  $L:\R^N\to\R^{4a+b}$ is linear.  This explicit parameter records every
  ignored-coordinate extension used below.
\end{definition}

\begin{lemma}[Depth transfer along $\Tspace_{a,b}(N)$]
  \label{lem:t-space-depth-transfer}
  \uses{def:t-space, lem:affine-precomposition,
    lem:parallel-linear-combination}
  If $T_{a,b}$ is computable with at most $k$ hidden layers, then every
  function in $\Tspace_{a,b}(N)$ is computable with at most $k$ hidden layers
  on $\R^N$.
\end{lemma}
\begin{proof}
  \uses{def:t-space, lem:affine-precomposition,
    lem:parallel-linear-combination}
  Every generator $T_{a,b}\circ L$ has depth at most $k$ by affine
  precomposition.  Every vector of the span is a finite linear combination of
  such generators, and the parallel linear-combination lemma preserves the
  same depth.
\end{proof}

\begin{lemma}[Lifted $P$-terms in Claim 5]
  \label{lem:claim-five-p-membership}
  \uses{def:t-function, def:t-space}
  Put $i=4a+b$ and replace $x_5$ in $P_1,P_2,P_3,P_4$ by
  $T_{a,b}(x_1,\ldots,x_i)$ while replacing $x_1,\ldots,x_4$ by the four new
  coordinates $x_{i+1},\ldots,x_{i+4}$.  Each resulting function belongs to
  $\Tspace_{a+1,b+1}(4a+b+5)$.
\end{lemma}
\begin{proof}
  \uses{def:t-function, def:t-space}
  For example the lifted $P_1$ is
  \[
  2\max\left(T_{a,b},\frac{x_{i+1}+x_{i+2}}2,
       \max\!\left(\frac{x_{i+1}}2,\frac{x_{i+3}}2\right)+
       \max\!\left(\frac{x_{i+1}}2,\frac{x_{i+4}}2\right)\right).
  \]
  A linear coordinate map feeds the old $4a+b$ coordinates into the
  $T_{a,b}$ part of $T_{a+1,b+1}$, feeds the indicated halved new coordinates
  into its one additional four-coordinate block, and uses one unused extra
  coordinate.  Multiplication by $2$ is allowed in the span.  The same map
  with the evident coordinate permutations gives $P_2,P_3,P_4$.
\end{proof}

\begin{lemma}[Lifted $Q$-term in Claim 5]
  \label{lem:claim-five-q-membership}
  \uses{def:t-function, def:t-space}
  Under the substitution of Lemma~\ref{lem:claim-five-p-membership}, the
  resulting $Q$ belongs to $\Tspace_{a,b+2}(4a+b+5)$.
\end{lemma}
\begin{proof}
  \uses{def:t-function, def:t-space}
  The lifted term is
  \[
    2\max\left(T_{a,b},\frac{x_{i+1}+x_{i+2}}2,
                         \frac{x_{i+3}+x_{i+4}}2\right).
  \]
  A linear map supplies the two displayed half-sums as the two extra scalar
  arguments of $T_{a,b+2}$; scaling by $2$ is allowed in its span.
\end{proof}

\begin{lemma}[Lifted $R$-terms in Claim 5]
  \label{lem:claim-five-r-membership}
  \uses{def:t-function, def:t-space}
  Under the same substitution, each of the four resulting terms
  $R_{13},R_{14},R_{23},R_{24}$ belongs to
  $\Tspace_{a,b+3}(4a+b+5)$.
\end{lemma}
\begin{proof}
  \uses{def:t-function, def:t-space}
  For example, the lifted $R_{13}$ is
  \[
    2\max\left(T_{a,b},\frac{x_{i+1}+x_{i+3}}2,
      \frac{x_{i+1}+x_{i+2}}2,\frac{x_{i+3}+x_{i+4}}2\right).
  \]
  Feed the last three linear half-sums to the three extra scalar slots of
  $T_{a,b+3}$ and multiply by $2$.  Changing only the middle half-sum gives
  the other three indexed terms.
\end{proof}

\begin{lemma}[First padding containment]
  \label{lem:tspace-b2-b3}
  \uses{def:t-space}
  For every ambient $N$ and all $a,b\geq0$,
  $\Tspace_{a,b+2}(N)\subseteq\Tspace_{a,b+3}(N)$.
\end{lemma}
\begin{proof}
  \uses{def:t-function, def:t-space}
  Extend every generator $T_{a,b+2}\circ L$ by ignoring one final coordinate.
  In the corresponding generator of $\Tspace_{a,b+3}(N)$, use a linear map that
  duplicates one of the existing scalar arguments into the added slot.  A
  duplicate does not change a maximum.  Linear combinations give the stated
  containment.
\end{proof}

\begin{lemma}[Second padding containment]
  \label{lem:tspace-b3-next}
  \uses{def:t-space}
  For every ambient $N$ and all $a,b\geq0$,
  $\Tspace_{a,b+3}(N)\subseteq\Tspace_{a+1,b+1}(N)$.
\end{lemma}
\begin{proof}
  \uses{def:t-function, def:t-space}
  The identity
  \[
    \max(u,v)=\max\bigl(\max(u,v)+\max(0,0)\bigr)
  \]
  represents two scalar arguments of $T_{a,b+3}$ by the one new four-input
  block of $T_{a+1,b+1}$: feed $(u,v,0,0)$ to that block.  Feed the remaining
  $b+1$ scalar arguments to the remaining scalar slots.  This realizes each
  generator by linear precomposition, and hence realizes its span.
\end{proof}

\begin{lemma}[Claim 5: the four-to-one span step]
  \label{lem:claim-five}
  \uses{def:t-space}
  Viewing $T_{a,b+4}$ as a function on $\R^{4a+b+5}$ that ignores the last
  coordinate, one has
  \[
    T_{a,b+4}\circ\pi\in\Tspace_{a+1,b+1}(4a+b+5),
  \]
  where $\pi:\R^{4a+b+5}\to\R^{4a+b+4}$ forgets the last coordinate.
\end{lemma}
\begin{proof}
  \uses{lem:claim-four, lem:claim-five-p-membership,
    lem:claim-five-q-membership, lem:claim-five-r-membership,
    lem:tspace-b2-b3, lem:tspace-b3-next}
  Write $i=4a+b$.  Then
  $T_{a,b+4}=\max(x_{i+1},x_{i+2},x_{i+3},x_{i+4},T_{a,b})$.
  Apply the signed identity of Claim~\ref{lem:claim-four} to these five
  arguments.  The four $P$-terms already lie in
  $\Tspace_{a+1,b+1}(4a+b+5)$.  The $Q$-term and the four $R$-terms lie in the
  corresponding same-ambient spaces from the preceding lemmas; the two
  padding containments put all of them in the target space.  Since this is a
  vector space, their signed half-sum is
  also in it.
\end{proof}

\begin{lemma}[Iterated Claim 5 containment]
  \label{lem:claim-five-iterated}
  \uses{lem:claim-five}
  For every $a,b,t\geq0$ with $b\geq3t$, let
  $\pi_t:\R^{4a+b+t}\to\R^{4a+b}$ forget the last $t$ coordinates.  Then
  \[
    T_{a,b}\circ\pi_t\in\Tspace_{a+t,b-3t}(4a+b+t).
  \]
\end{lemma}
\begin{proof}
  \uses{lem:claim-five, def:t-space}
  Induct on $t$.  The case $t=0$ is the defining generator with the identity
  linear map.  For the successor, write the current scalar count as
  $(b-3t-4)+4$ and apply Claim~\ref{lem:claim-five}, then use linearity and
  precomposition to transport every element of the current span into the next
  span.  The indices change from $(a+t,b-3t)$ to
  $(a+t+1,b-3(t+1))$, as required.
\end{proof}

\begin{lemma}[Claim 6: transfer from blocks to a large maximum]
  \label{lem:claim-six}
  \uses{def:t-function}
  Let $r\geq1$.  If $T_{3^{r-1},2}$ is computable with at most $k$ hidden
  layers, then $T_{0,3^r+2}=\MAX_{3^r+2}$ is computable with at most $k$
  hidden layers.
\end{lemma}
\begin{proof}
  \uses{lem:claim-five-iterated, lem:t-space-depth-transfer}
  Apply Lemma~\ref{lem:claim-five-iterated} with
  $(a,b,t)=(0,3^r+2,3^{r-1})$.  Since
  $3^r+2-3\cdot3^{r-1}=2$, it gives
  the ignored-coordinate extension of $T_{0,3^r+2}$ in
  $\Tspace_{3^{r-1},2}(3^r+2+3^{r-1})$.  The depth-transfer lemma now gives
  the same $k$-layer bound as for $T_{3^{r-1},2}$.
\end{proof}

\begin{lemma}[Claim 7: one layer creates all four-input blocks]
  \label{lem:claim-seven}
  \uses{def:t-function, lem:max-two-relu}
  Let $r\geq1$.  If $T_{0,3^{r-1}+2}$ is computable with at most $k-1$
  hidden layers, then $T_{3^{r-1},2}$ is computable with at most $k$ hidden
  layers.
\end{lemma}
\begin{proof}
  \uses{lem:max-two-relu, def:t-function}
  In one hidden layer, for each of the $3^{r-1}$ disjoint four-input blocks,
  compute both pairwise maxima and add them affinely, yielding
  $\max(x_{4j+1},x_{4j+2})+\max(x_{4j+3},x_{4j+4})$.
  Carry the two extra inputs through that layer using positive/negative ReLU
  pairs.  Apply the assumed realization of $T_{0,3^{r-1}+2}$ to the
  $3^{r-1}$ block values and those two inputs.  This uses one plus $k-1$
  hidden layers and computes exactly $T_{3^{r-1},2}$.
\end{proof}

\begin{lemma}[Base-three ceiling arithmetic]
  \label{lem:ceiling-log-three}
  For every integer $m\geq4$, if
  $r=\lceil\log_3(m-2)\rceil$, then $r\geq1$ and $m\leq3^r+2$.
\end{lemma}
\begin{proof}
  Since $m-2\geq2>1$, its base-three logarithm is positive, hence $r\geq1$.
  By the defining property of the ceiling,
  $\log_3(m-2)\leq r$.  The function $t\mapsto3^t$ is increasing, so
  $m-2\leq3^r$, and adding $2$ proves the inequality.
\end{proof}

\begin{corollary}[Maximum bound for arbitrary arity]
  \label{cor:max-general-depth}
  \uses{thm:main-max, lem:ceiling-log-three, lem:max-arity-restriction}
  For every integer $m\geq4$,
  \[
    \MAX_m\in\ReLU_{m,\lceil\log_3(m-2)\rceil+1}.
  \]
\end{corollary}
\begin{proof}
  \uses{thm:main-max, lem:ceiling-log-three, lem:max-arity-restriction}
  Set $r=\lceil\log_3(m-2)\rceil$.  The main maximum theorem realizes
  $\MAX_{3^r+2}$ with $r+1$ layers, and
  Lemma~\ref{lem:ceiling-log-three} gives $m\leq3^r+2$.  Restrict the larger
  maximum to arity $m$.
\end{proof}

\begin{corollary}[Four-dimensional CPWL functions need at most two layers]
  \label{cor:cpwl-four-two-layers}
  \uses{thm:main-cpwl}
  One has $\CPWL_4=\ReLU_{4,2}$.
\end{corollary}
\begin{proof}
  \uses{thm:main-cpwl}
  Substitute $n=4$ in Theorem~\ref{thm:main-cpwl}; since
  $\lceil\log_3(3)\rceil+1=2$, the displayed equality is exactly the result.
\end{proof}

\begin{proposition}[Remark 1: dyadic weights]
  \label{prop:dyadic-weights}
  \uses{lem:claim-five, lem:claim-seven, lem:max-five-upper-bound}
  Every network produced by the induction proving Theorem~\ref{thm:main-max}
  can be chosen with weights whose denominators are powers of two.
\end{proposition}
\begin{proof}
  \uses{def:m-expression, lem:claim-five-p-membership,
    lem:claim-five-q-membership, lem:claim-five-r-membership,
    lem:claim-seven}
  The base construction uses integer affine forms and the outer coefficient
  $1/2$.  In the Claim 5 step, every new coordinate coefficient is obtained
  from an old coefficient by integer addition, sign change, or division by
  $2$, as the three explicit membership formulas show.  Claim 7 introduces
  only the integer coefficients occurring in pairwise maxima and sums.
  Induction therefore keeps all weights in $\mathbb Z[1/2]$, precisely the
  fractions with power-of-two denominators.
\end{proof}

\begin{corollary}[Dyadic weights are decimal fractions]
  \label{cor:dyadic-decimal}
  \uses{prop:dyadic-weights}
  All weights in Proposition~\ref{prop:dyadic-weights} have a finite decimal
  expansion.
\end{corollary}
\begin{proof}
  A weight $a/2^q$ equals $a5^q/10^q$, which is a decimal fraction.
\end{proof}

\begin{proposition}[Linear neuron count]
  \label{prop:linear-size}
  \uses{def:relu-network, cor:max-general-depth}
  \notready
  There is an absolute constant $C$ such that the paper's construction of
  $\MAX_m$ uses at most $Cm$ hidden neurons for every $m\geq4$.
\end{proposition}
% TODO: The paper says this is easy to verify but supplies neither an explicit
% recurrence nor a constant.  Record exact widths for Claims 5 and 7 first.

\begin{definition}[The depth classes $\CPWL_n$ and $\ReLU_{n,k}$]
  \label{def:depth-classes}
  \uses{def:relu-network, def:cpwl-function, lem:relu-network-is-cpwl}
  Write $\CPWL_n$ for all CPWL functions $\R^n\to\R$, and write
  $\ReLU_{n,k}\subseteq\CPWL_n$ for those functions computable by a ReLU
  network with at most $k$ hidden layers.
\end{definition}

\begin{definition}[Maximum function]
  \label{def:max-function}
  For $m\geq1$, define
  \[
    \MAX_m(x_1,\ldots,x_m)=\max\{x_1,\ldots,x_m\}.
  \]
\end{definition}

\begin{theorem}[Wang--Sun signed maximum representation, Equation (1)]
  \label{thm:wang-sun-representation}
  \uses{def:cpwl-function, def:max-function}
  \notready
  If $f\in\CPWL_n$, then there are $s\geq1$, affine maps
  $A_i:\R^n\to\R^{n+1}$, and signs $\sigma_i\in\{-1,1\}$ such that
  \[
    f(x)=\sum_{i=1}^s\sigma_i\MAX_{n+1}(A_i(x))
    \qquad(x\in\R^n).
  \]
\end{theorem}
% TODO: The paper cites [WS05] and does not reproduce the proof.  Import or
% formalize the full hinging-hyperplane argument before marking this ready.

\begin{lemma}[Pairwise maximum is one ReLU]
  \label{lem:max-two-relu}
  \uses{def:relu-scalar, def:relu-network, def:max-function}
  For all $u,v\in\R$,
  \[
    \max\{u,v\}=u+\ReLU(v-u).
  \]
  Consequently $\MAX_2$ is computable with one hidden layer.
\end{lemma}
\begin{proof}
  \uses{def:relu-scalar, def:relu-network}
  If $v\leq u$, then $v-u\leq0$, so the right side is $u$.  If $u\leq v$,
  then $\ReLU(v-u)=v-u$, so the right side is $v$.  These cases cover all
  pairs.  The displayed formula is an affine output applied to the single
  hidden neuron $\ReLU(v-u)$, with the affine input $u$ carried to the output;
  the standard identity $u=\ReLU(u)-\ReLU(-u)$ supplies an ordinary
  layer-by-layer realization if direct affine skip terms are disallowed.
\end{proof}

\begin{lemma}[Affine precomposition preserves depth]
  \label{lem:affine-precomposition}
  \uses{def:relu-network}
  If $g:\R^q\to\R$ has a realization with at most $k$ hidden layers and
  $A:\R^p\to\R^q$ is affine, then $g\circ A$ has a realization with at most
  $k$ hidden layers.
\end{lemma}
\begin{proof}
  \uses{def:relu-network}
  In a realization of $g$, replace its first affine transformation $T^{(1)}$
  by $T^{(1)}\circ A$.  This is affine, all later transformations are
  unchanged, and the resulting alternating composition is exactly $g\circ A$.
\end{proof}

\begin{lemma}[Parallel linear combinations preserve depth]
  \label{lem:parallel-linear-combination}
  \uses{def:relu-network}
  Let $g_1,\ldots,g_s:\R^n\to\R$ each be computable with at most $k$ hidden
  layers.  For arbitrary real $c_1,\ldots,c_s$, the function
  $\sum_i c_i g_i$ is computable with at most $k$ hidden layers.
\end{lemma}
\begin{proof}
  \uses{def:relu-network}
  Pad a shallower realization by identity layers, implementing each scalar
  identity as $z=\ReLU(z)-\ReLU(-z)$, until every branch has depth $k$.
  At every hidden layer take the block-diagonal direct sum of the affine maps,
  so all networks run in parallel.  The coordinatewise ReLU respects these
  blocks.  In the final affine map take the linear combination of the branch
  outputs with coefficients $c_i$.  No additional hidden layer is introduced.
\end{proof}

\begin{lemma}[Restriction from a larger maximum]
  \label{lem:max-arity-restriction}
  \uses{def:max-function, lem:affine-precomposition}
  If $1\leq m\leq N$ and $\MAX_N$ is computable with $k$ hidden layers, then
  $\MAX_m$ is computable with $k$ hidden layers.
\end{lemma}
\begin{proof}
  \uses{lem:affine-precomposition}
  Define the affine map
  $A(x_1,\ldots,x_m)=(x_1,\ldots,x_m,x_1,\ldots,x_1)\in\R^N$.
  Its last $N-m$ coordinates repeat $x_1$, so
  $\MAX_N\circ A=\MAX_m$.  Apply Lemma~\ref{lem:affine-precomposition}.
\end{proof}

\begin{proposition}[Binary-tree maximum bound]
  \label{prop:binary-tree-max-bound}
  \uses{lem:max-two-relu, def:max-function}
  For every $m\geq1$, $\MAX_m$ is computable with at most
  $\lceil\log_2m\rceil$ hidden layers.
\end{proposition}
\begin{proof}
  \uses{lem:max-two-relu, lem:max-arity-restriction}
  Let $d=\lceil\log_2m\rceil$ and repeat $x_1$ to regard the input as a list of
  $2^d$ numbers without changing its maximum.  In the first hidden layer use
  Lemma~\ref{lem:max-two-relu} in parallel on adjacent pairs.  Repeat on the
  $2^{d-1}$ outputs, then on the $2^{d-2}$ outputs, and so on.  After exactly
  $d$ pairwise-max layers one number remains, and associativity of maximum says
  that it is the maximum of the padded list, hence $\MAX_m$.
\end{proof}

\begin{theorem}[Previous CPWL depth bound]
  \label{thm:previous-cpwl-bound}
  \uses{def:depth-classes, thm:wang-sun-representation,
    prop:binary-tree-max-bound}
  For every $n\geq1$,
  \[
    \CPWL_n\subseteq\ReLU_{n,\lceil\log_2(n+1)\rceil}.
  \]
\end{theorem}
\begin{proof}
  \uses{thm:wang-sun-representation, prop:binary-tree-max-bound,
    lem:affine-precomposition, lem:parallel-linear-combination}
  Write $f$ in the form of Theorem~\ref{thm:wang-sun-representation}.
  Proposition~\ref{prop:binary-tree-max-bound} realizes $\MAX_{n+1}$ with the
  displayed depth.  Precompose each copy by $A_i$ and run all signed copies in
  parallel.  Their affine output sum computes $f$ at the same depth.
\end{proof}

\begin{theorem}[Mukherjee--Basu lower bound for $\MAX_3$]
  \label{thm:max-three-lower-bound}
  \uses{def:relu-network, def:max-function}
  \notready
  No ReLU network with only one hidden layer computes $\MAX_3$ on all of
  $\R^3$.
\end{theorem}
% TODO: The paper cites [MB17] for this lower bound and omits its proof.

\begin{theorem}[Theorem 1: base-three maximum bound]
  \label{thm:main-max}
  \uses{def:depth-classes, def:max-function}
  For every integer $r\geq1$,
  \[
    \MAX_{3^r+2}\in\ReLU_{3^r+2,r+1}.
  \]
\end{theorem}
\begin{proof}
  \uses{lem:max-five-upper-bound, lem:claim-six, lem:claim-seven,
    def:t-function}
  Induct on $r$.  For $r=1$, $3^r+2=5$, and
  Lemma~\ref{lem:max-five-upper-bound} gives two hidden layers.  For the
  induction step, assume
  $T_{0,3^{r-1}+2}=\MAX_{3^{r-1}+2}$ is computable with $r$ hidden layers.
  Claim~\ref{lem:claim-seven} computes $T_{3^{r-1},2}$ with $r+1$ hidden
  layers.  Claim~\ref{lem:claim-six} then computes
  $T_{0,3^r+2}=\MAX_{3^r+2}$ with those same $r+1$ layers.
\end{proof}

\begin{theorem}[Theorem 2: improved depth for every CPWL function]
  \label{thm:main-cpwl}
  \uses{def:depth-classes}
  For every $n\geq3$,
  \[
    \CPWL_n=\ReLU_{n,\lceil\log_3(n-1)\rceil+1}.
  \]
\end{theorem}
\begin{proof}
  \uses{thm:wang-sun-representation, cor:max-general-depth,
    lem:affine-precomposition, lem:parallel-linear-combination,
    lem:relu-network-is-cpwl}
  Let $f\in\CPWL_n$.  The Wang--Sun representation writes $f$ as a signed
  sum of affine precompositions of $\MAX_{n+1}$.  Since $n\geq3$, the general
  maximum bound realizes $\MAX_{n+1}$ with
  $\lceil\log_3((n+1)-2)\rceil+1=
  \lceil\log_3(n-1)\rceil+1$ hidden layers.  Affine precomposition and
  parallel signed summation preserve this depth, proving
  $\CPWL_n\subseteq\ReLU_{n,\lceil\log_3(n-1)\rceil+1}$.  The reverse
  inclusion is Lemma~\ref{lem:relu-network-is-cpwl} together with the
  definition of the depth class.
\end{proof}

\chapter{Direct proofs}

\section{Computing the maximum of five inputs}

\begin{definition}[The term $P_1$ in Equation (2)]
  \label{def:p1}
  Define $P_1:\R^5\to\R$ by
  \[
  P_1(x)=\max\!\bigl(\max(2x_5,x_1+x_2),
               \max(x_1,x_3)+\max(x_1,x_4)\bigr).
  \]
  Equivalently it is the maximum of
  $2x_5,x_1+x_2,2x_1,x_1+x_3,x_1+x_4,x_3+x_4$.
\end{definition}

\begin{definition}[The term $P_2$ in Equation (2)]
  \label{def:p2}
  Define
  \[
  P_2(x)=\max\!\bigl(\max(2x_5,x_1+x_2),
               \max(x_2,x_3)+\max(x_2,x_4)\bigr),
  \]
  equivalently the maximum of
  $2x_5,x_1+x_2,2x_2,x_2+x_3,x_2+x_4,x_3+x_4$.
\end{definition}

\begin{definition}[The term $P_3$ in Equation (2)]
  \label{def:p3}
  Define
  \[
  P_3(x)=\max\!\bigl(\max(2x_5,x_3+x_4),
               \max(x_3,x_1)+\max(x_3,x_2)\bigr),
  \]
  equivalently the maximum of
  $2x_5,x_3+x_4,2x_3,x_3+x_1,x_3+x_2,x_1+x_2$.
\end{definition}

\begin{definition}[The term $P_4$ in Equation (2)]
  \label{def:p4}
  Define
  \[
  P_4(x)=\max\!\bigl(\max(2x_5,x_3+x_4),
               \max(x_4,x_1)+\max(x_4,x_2)\bigr),
  \]
  equivalently the maximum of
  $2x_5,x_3+x_4,2x_4,x_4+x_1,x_4+x_2,x_1+x_2$.
\end{definition}

\begin{definition}[The term $Q$ in Equation (2)]
  \label{def:q}
  Define
  \[
    Q(x)=\max\bigl(2x_5,\max(x_1+x_2,x_3+x_4)\bigr)
        =\max(2x_5,x_1+x_2,x_3+x_4).
  \]
\end{definition}

\begin{definition}[The term $R_{13}$ in Equation (2)]
  \label{def:r13}
  Define
  \[
  R_{13}(x)=\max\bigl(\max(2x_5,x_1+x_3),
                       \max(x_1+x_2,x_3+x_4)\bigr).
  \]
\end{definition}

\begin{definition}[The term $R_{14}$ in Equation (2)]
  \label{def:r14}
  Define
  \[
  R_{14}(x)=\max\bigl(\max(2x_5,x_1+x_4),
                       \max(x_1+x_2,x_3+x_4)\bigr).
  \]
\end{definition}

\begin{definition}[The term $R_{23}$ in Equation (2)]
  \label{def:r23}
  Define
  \[
  R_{23}(x)=\max\bigl(\max(2x_5,x_2+x_3),
                       \max(x_1+x_2,x_3+x_4)\bigr).
  \]
\end{definition}

\begin{definition}[The term $R_{24}$ in Equation (2)]
  \label{def:r24}
  Define
  \[
  R_{24}(x)=\max\bigl(\max(2x_5,x_2+x_4),
                       \max(x_1+x_2,x_3+x_4)\bigr).
  \]
\end{definition}

\begin{definition}[The signed expression $M$, Equation (2)]
  \label{def:m-expression}
  \uses{def:p1, def:p2, def:p3, def:p4, def:q, def:r13, def:r14,
    def:r23, def:r24}
  Define
  \[
  M(x)=\frac12\bigl(P_1(x)+P_2(x)+P_3(x)+P_4(x)+Q(x)
             -R_{13}(x)-R_{14}(x)-R_{23}(x)-R_{24}(x)\bigr).
  \]
\end{definition}

\begin{lemma}[Two-layer max-of-maxes template]
  \label{lem:two-layer-template}
  \uses{def:relu-network, lem:max-two-relu}
  If $a,b,c,d,e,f:\R^n\to\R$ are affine, then
  \[
    x\longmapsto
    \max\bigl(\max(a(x),b(x)),\max(c(x),d(x))+\max(e(x),f(x))\bigr)
  \]
  is computable with at most two hidden layers.
\end{lemma}
\begin{proof}
  \uses{lem:max-two-relu, lem:parallel-linear-combination}
  In the first hidden layer compute the three pairwise maxima
  $u=\max(a,b)$, $v=\max(c,d)$, and $w=\max(e,f)$ using
  Lemma~\ref{lem:max-two-relu} in parallel.  Their sum $v+w$ is affine in the
  first-layer outputs.  In the second hidden layer compute
  $\max(u,v+w)$ by the same identity.  This is the required function.
\end{proof}

\begin{lemma}[Each $P_i$ has depth two]
  \label{lem:p-terms-depth-two}
  \uses{def:p1, def:p2, def:p3, def:p4, lem:two-layer-template}
  Each of $P_1,P_2,P_3,P_4$ is computable with at most two hidden layers.
\end{lemma}
\begin{proof}
  \uses{lem:two-layer-template}
  For $P_1$, substitute
  $(a,b,c,d,e,f)=(2x_5,x_1+x_2,x_1,x_3,x_1,x_4)$ in
  Lemma~\ref{lem:two-layer-template}.  Replacing the repeated coordinate
  $x_1$ successively by $x_2$, and then using the pair $(x_3,x_4)$ in the
  first inner maximum and repeating $x_3$ or $x_4$, gives exactly the displayed
  definitions of $P_2,P_3,P_4$.
\end{proof}

\begin{lemma}[$Q$ has depth two]
  \label{lem:q-depth-two}
  \uses{def:q, lem:two-layer-template}
  The function $Q$ is computable with at most two hidden layers.
\end{lemma}
\begin{proof}
  \uses{lem:two-layer-template}
  Apply the template with
  $(a,b,c,d,e,f)=(2x_5,2x_5,x_1+x_2,x_3+x_4,0,0)$.
  Then the first pairwise maximum is $2x_5$, the last is zero, and the result
  is $\max(2x_5,\max(x_1+x_2,x_3+x_4))=Q$.
\end{proof}

\begin{lemma}[Each $R_{ij}$ has depth two]
  \label{lem:r-terms-depth-two}
  \uses{def:r13, def:r14, def:r23, def:r24, lem:two-layer-template}
  Each of $R_{13},R_{14},R_{23},R_{24}$ is computable with at most two hidden
  layers.
\end{lemma}
\begin{proof}
  \uses{lem:two-layer-template}
  For $R_{ij}$, apply the template with
  $(a,b)=(2x_5,x_i+x_j)$,
  $(c,d)=(x_1+x_2,x_3+x_4)$, and $(e,f)=(0,0)$.
  The final pairwise maximum contributes zero, leaving exactly the definition
  of $R_{ij}$ for each of the four indicated pairs.
\end{proof}

\begin{lemma}[$x_5$-dominant case of Claim 4]
  \label{lem:m-x5-case}
  \uses{def:m-expression, def:max-function}
  If $x_5=\MAX_5(x)$, then $M(x)=x_5$.
\end{lemma}
\begin{proof}
  \uses{def:p1, def:p2, def:p3, def:p4, def:q, def:r13, def:r14,
    def:r23, def:r24}
  For $i,j\in\{1,2,3,4\}$, the hypotheses $x_i\leq x_5$ and
  $x_j\leq x_5$ give $x_i+x_j\leq2x_5$.  Thus every displayed candidate in
  each $P_i,Q,R_{ij}$ is at most $2x_5$, while $2x_5$ itself is a candidate.
  All nine terms equal $2x_5$.  Equation (2) has five positive and four
  negative such terms, so $M(x)=\tfrac12(2x_5)=x_5$.
\end{proof}

\begin{lemma}[Symmetries of $M$]
  \label{lem:m-symmetry}
  \uses{def:m-expression}
  The function $M$ is invariant under swapping $1$ with $2$, under swapping
  $3$ with $4$, and under simultaneously swapping the pairs $(1,2)$ and
  $(3,4)$.
\end{lemma}
\begin{proof}
  \uses{def:p1, def:p2, def:p3, def:p4, def:q, def:r13, def:r14,
    def:r23, def:r24}
  Under $1\leftrightarrow2$, the terms $P_1,P_2$ are exchanged,
  $R_{13},R_{23}$ are exchanged, and $R_{14},R_{24}$ are exchanged; all other
  terms are fixed.  Under $3\leftrightarrow4$, the analogous exchanges are
  $P_3\leftrightarrow P_4$, $R_{13}\leftrightarrow R_{14}$, and
  $R_{23}\leftrightarrow R_{24}$.  Under the pair swap,
  $P_1\leftrightarrow P_3$, $P_2\leftrightarrow P_4$, and
  $R_{14}\leftrightarrow R_{23}$, with all other signed terms fixed.
  In every case the signed sum defining $M$ is unchanged.
\end{proof}

\begin{lemma}[$P_1$ in the $x_1$-dominant case]
  \label{lem:p1-x1-case}
  \uses{def:p1}
  If $x_1\geq x_j$ for $j=2,3,4,5$, then $P_1(x)=2x_1$.
\end{lemma}
\begin{proof}
  \uses{def:p1}
  Each candidate other than $2x_1$ is a sum of two coordinates, both at most
  $x_1$, or is $2x_5\leq2x_1$.  Hence every candidate is at most $2x_1$, and
  $2x_1$ is itself present in the maximum.
\end{proof}

\begin{lemma}[$P_2$ cancels $R_{23}$]
  \label{lem:p2-r23-x1-case}
  \uses{def:p2, def:r23}
  If $x_1\geq x_j$ for $j=2,3,4,5$, then $P_2(x)=R_{23}(x)$.
\end{lemma}
\begin{proof}
  \uses{def:p2, def:r23}
  We have $x_1+x_2\geq2x_2$ and
  $x_1+x_2\geq x_2+x_4$.  Removing these dominated candidates from the
  maximum defining $P_2$ leaves
  $\max(2x_5,x_1+x_2,x_2+x_3,x_3+x_4)$, which is $R_{23}$.
\end{proof}

\begin{lemma}[$P_3$ cancels $R_{13}$]
  \label{lem:p3-r13-x1-case}
  \uses{def:p3, def:r13}
  If $x_1\geq x_j$ for $j=2,3,4,5$, then $P_3(x)=R_{13}(x)$.
\end{lemma}
\begin{proof}
  \uses{def:p3, def:r13}
  The inequalities $x_1+x_3\geq2x_3$ and
  $x_1+x_3\geq x_2+x_3$ show that those two candidates are redundant in
  $P_3$.  The remaining maximum is
  $\max(2x_5,x_3+x_4,x_1+x_3,x_1+x_2)=R_{13}$.
\end{proof}

\begin{lemma}[$P_4$ cancels $R_{14}$]
  \label{lem:p4-r14-x1-case}
  \uses{def:p4, def:r14}
  If $x_1\geq x_j$ for $j=2,3,4,5$, then $P_4(x)=R_{14}(x)$.
\end{lemma}
\begin{proof}
  \uses{def:p4, def:r14}
  The inequalities $x_1+x_4\geq2x_4$ and
  $x_1+x_4\geq x_2+x_4$ remove those dominated candidates from $P_4$.
  What remains is
  $\max(2x_5,x_3+x_4,x_1+x_4,x_1+x_2)=R_{14}$.
\end{proof}

\begin{lemma}[$R_{24}$ cancels $Q$]
  \label{lem:r24-q-x1-case}
  \uses{def:r24, def:q}
  If $x_1\geq x_j$ for $j=2,3,4,5$, then $R_{24}(x)=Q(x)$.
\end{lemma}
\begin{proof}
  \uses{def:r24, def:q}
  Since $x_1+x_2\geq x_2+x_4$, the candidate $x_2+x_4$ is redundant in
  $R_{24}$.  Its remaining three candidates are precisely
  $2x_5,x_1+x_2,x_3+x_4$, the candidates defining $Q$.
\end{proof}

\begin{lemma}[$x_1$-dominant case of Claim 4]
  \label{lem:m-x1-case}
  \uses{def:m-expression, lem:p1-x1-case, lem:p2-r23-x1-case,
    lem:p3-r13-x1-case, lem:p4-r14-x1-case, lem:r24-q-x1-case}
  If $x_1\geq x_j$ for $j=2,3,4,5$, then $M(x)=x_1$.
\end{lemma}
\begin{proof}
  \uses{lem:p1-x1-case, lem:p2-r23-x1-case, lem:p3-r13-x1-case,
    lem:p4-r14-x1-case, lem:r24-q-x1-case}
  In Equation (2), $P_2$ cancels $R_{23}$, $P_3$ cancels $R_{13}$,
  $P_4$ cancels $R_{14}$, and $Q$ cancels $R_{24}$.  Thus
  $M(x)=\tfrac12P_1(x)=\tfrac12(2x_1)=x_1$.
\end{proof}

\begin{lemma}[Claim 4: the signed identity]
  \label{lem:claim-four}
  \uses{def:m-expression, def:max-function}
  For every $x\in\R^5$, $M(x)=\MAX_5(x)$.
\end{lemma}
\begin{proof}
  \uses{lem:m-x5-case, lem:m-symmetry, lem:m-x1-case}
  If $x_5$ is maximal, use Lemma~\ref{lem:m-x5-case}.  Otherwise one of
  $x_1,x_2,x_3,x_4$ is maximal.  The three symmetries in
  Lemma~\ref{lem:m-symmetry} act transitively on these four indices, so a
  symmetry sends a maximal coordinate to position $1$.  The value of $M$ is
  unchanged, and Lemma~\ref{lem:m-x1-case} gives that value as the maximal
  coordinate.  The ordinary maximum is also invariant under the same
  permutation, proving the identity.
\end{proof}

\begin{lemma}[Two-layer upper bound for $\MAX_5$]
  \label{lem:max-five-upper-bound}
  \uses{lem:claim-four, lem:p-terms-depth-two, lem:q-depth-two,
    lem:r-terms-depth-two}
  The function $\MAX_5$ is computable with at most two hidden layers.
\end{lemma}
\begin{proof}
  \uses{def:m-expression, lem:parallel-linear-combination,
    lem:p-terms-depth-two, lem:q-depth-two, lem:r-terms-depth-two,
    lem:claim-four}
  Compute the nine terms of Equation (2) in parallel at depth two, and use the
  affine output layer to form their signed half-sum.  By
  Lemma~\ref{lem:claim-four}, this output is $\MAX_5$.
\end{proof}

\begin{proposition}[Proposition 3: exact depth of $\MAX_5$]
  \label{prop:max-five-depth}
  \uses{lem:max-five-upper-bound, thm:max-three-lower-bound,
    lem:max-arity-restriction}
  The minimum number of hidden layers required to compute $\MAX_5$ is exactly
  two.
\end{proposition}
\begin{proof}
  \uses{lem:max-five-upper-bound, thm:max-three-lower-bound,
    lem:max-arity-restriction}
  The upper bound is Lemma~\ref{lem:max-five-upper-bound}.  If one hidden layer
  computed $\MAX_5$, Lemma~\ref{lem:max-arity-restriction} with $m=3$ would
  give a one-hidden-layer realization of $\MAX_3$, contradicting
  Theorem~\ref{thm:max-three-lower-bound}.  Thus two layers are also necessary.
\end{proof}

\chapter{Geometric intuition}

\section{Polytopes, support functions, and neural networks}

\begin{definition}[Polytope]
  \label{def:polytope}
  A polytope in $\R^n$ is the convex hull of a nonempty finite set of points.
\end{definition}

\begin{definition}[Support function]
  \label{def:support-function}
  \uses{def:polytope}
  For a polytope $P\subseteq\R^n$, its support function is
  \[
    h_P(x)=\max\{\langle x,p\rangle:p\in P\}.
  \]
\end{definition}

\begin{lemma}[Support functions of polytopes]
  \label{lem:support-function-properties}
  \uses{def:support-function, def:cpwl-function}
  The support function of a polytope is convex, positively homogeneous, and
  CPWL.
\end{lemma}
\begin{proof}
  \uses{def:polytope, def:support-function}
  Write $P=\operatorname{conv}\{v_1,\ldots,v_s\}$.  A linear functional
  attains its maximum over this convex hull at one of the listed vertices, so
  $h_P(x)=\max_i\langle x,v_i\rangle$.  This finite maximum of linear maps is
  positively homogeneous.  It is convex because for $0\leq t\leq1$,
  \[
  \max_i\langle tx+(1-t)y,v_i\rangle
  \leq t\max_i\langle x,v_i\rangle+(1-t)\max_i\langle y,v_i\rangle.
  \]
  The finitely many regions on which a fixed linear form dominates are
  polyhedra; they cover $\R^n$, and the maximum is linear on each.  A finite
  maximum of continuous maps is continuous, hence the function is CPWL.
\end{proof}

\begin{theorem}[Polytope--support-function correspondence]
  \label{thm:support-polytope-correspondence}
  \uses{lem:support-function-properties}
  \notready
  Every convex, positively homogeneous CPWL function $f:\R^n\to\R$ is the
  support function of a unique polytope.
\end{theorem}
% TODO: This well-known correspondence is cited but not proved in the paper,
% and no matching theorem was found in the installed Mathlib/cslib sources.

\begin{definition}[Newton polytope]
  \label{def:newton-polytope}
  \uses{thm:support-polytope-correspondence}
  For a convex positively homogeneous CPWL function $f$, its Newton polytope
  is the unique polytope $P$ satisfying $f=h_P$.
\end{definition}

\begin{definition}[Mathlib standard simplex]
  \label{def:standard-simplex}
  \lean{stdSimplex}
  \mathlibok
  The standard simplex on a finite coordinate type is the set of nonnegative
  coordinate vectors whose coordinates sum to one.  In Mathlib it is the
  convex hull of the standard basis, by
  \texttt{convexHull\_basis\_eq\_stdSimplex}.
\end{definition}

\begin{lemma}[Newton polytope of the maximum]
  \label{lem:max-newton-simplex}
  \uses{def:newton-polytope, def:standard-simplex, def:max-function}
  The Newton polytope of $\MAX_n$ is
  $\Delta_{n-1}=\operatorname{conv}\{e_1,\ldots,e_n\}\subseteq\R^n$.
\end{lemma}
\begin{proof}
  \uses{def:support-function, thm:support-polytope-correspondence}
  A linear functional has the same maximum on a finite set and its convex
  hull.  Hence
  \[
    h_{\Delta_{n-1}}(x)=\max_i\langle x,e_i\rangle=\max_i x_i=\MAX_n(x).
  \]
  Uniqueness in the support-function correspondence identifies the Newton
  polytope with $\Delta_{n-1}$.
\end{proof}

\begin{definition}[Minkowski sum]
  \label{def:minkowski-sum}
  \uses{def:polytope}
  For polytopes $P,Q\subseteq\R^n$, define
  $P+Q=\{p+q:p\in P,\ q\in Q\}$.
\end{definition}

\begin{definition}[Mathlib convex join]
  \label{def:convex-join}
  \lean{convexJoin}
  \mathlibok
  The convex join of two sets is the union of all line segments joining one
  point of the first set to one point of the second.  For convex $P,Q$, it is
  $\operatorname{conv}(P\cup Q)$; the paper writes this polytope as $P*Q$.
\end{definition}

\begin{lemma}[Associativity of convex join]
  \label{lem:convex-join-assoc}
  \uses{def:convex-join}
  \lean{convexJoin_assoc}
  \mathlibok
  Convex join is associative: $(P*Q)*R=P*(Q*R)$.
\end{lemma}

\begin{definition}[The recursively generated classes $\PolyC_k$]
  \label{def:polytope-classes}
  \uses{def:polytope, def:minkowski-sum, def:convex-join}
  Fix an ambient dimension.  Let $\PolyC_0$ be the singleton polytopes.  For
  $k\geq0$, let $\PolyC_{k+1}$ consist of all finite Minkowski sums
  \[
    \sum_{i\in I}(P_i*Q_i),\qquad P_i,Q_i\in\PolyC_k,
  \]
  where $I$ is finite and nonempty.
\end{definition}

\begin{lemma}[Nesting of the polytope classes]
  \label{lem:polytope-class-monotone}
  \uses{def:polytope-classes}
  For every $k$, $\PolyC_k\subseteq\PolyC_{k+1}$.
\end{lemma}
\begin{proof}
  \uses{def:polytope-classes}
  If $P\in\PolyC_k$, then $P=P*P$, and this one-term sum is allowed in the
  definition of $\PolyC_{k+1}$.
\end{proof}

\begin{lemma}[Closure under finite Minkowski sums]
  \label{lem:polytope-class-sum-closed}
  \uses{def:polytope-classes}
  A finite Minkowski sum of members of $\PolyC_k$ again belongs to
  $\PolyC_k$.
\end{lemma}
\begin{proof}
  \uses{def:polytope-classes}
  Induct on $k$.  A finite sum of singleton polytopes is a singleton.  For the
  successor, expand every summand as a finite sum of joins of members of
  $\PolyC_k$.  Concatenating these finite indexing families expresses the
  total sum in exactly the defining form for $\PolyC_{k+1}$.
\end{proof}

\begin{definition}[Polytope complexity]
  \label{def:polytope-complexity}
  \uses{def:polytope-classes}
  The complexity of a polytope $P$ is the least $k$ such that
  $P\in\PolyC_k$, whenever such a $k$ exists.
\end{definition}

\begin{definition}[Formal Minkowski difference]
  \label{def:formal-minkowski-difference}
  \uses{def:minkowski-sum}
  A polytope $X$ is a formal Minkowski difference of $P$ and $Q$ if
  $X+P=Q$.
\end{definition}

\begin{lemma}[Lemma 8: ReLU depth as formal Minkowski difference]
  \label{lem:relu-depth-formal-difference}
  \uses{def:support-function, def:polytope-classes,
    def:formal-minkowski-difference, def:relu-network}
  \notready
  For every polytope $X$, the minimum hidden-layer depth representing $h_X$
  equals the least $k$ for which there are $P,Q\in\PolyC_k$ with $X+P=Q$.
\end{lemma}
% TODO: The paper cites [Her22, Theorem 3.35] and gives no proof.  Formalize
% the network/Newton-polytope semantics and both directions of that theorem.

\begin{corollary}[Maximum depth as a simplex difference]
  \label{cor:max-depth-formal-difference}
  \uses{lem:relu-depth-formal-difference, lem:max-newton-simplex}
  The minimum depth of $\MAX_n$ is the least $k$ for which
  $\Delta_{n-1}+P=Q$ for some $P,Q\in\PolyC_k$.
\end{corollary}
\begin{proof}
  \uses{lem:relu-depth-formal-difference, lem:max-newton-simplex}
  Substitute the Newton polytope $\Delta_{n-1}$ of $\MAX_n$ into
  Lemma~\ref{lem:relu-depth-formal-difference}.
\end{proof}

\section{Upper bounds from subdivisions}

\begin{lemma}[Two-cell valuation identity, Equation (3)]
  \label{lem:two-cell-valuation}
  \uses{def:minkowski-sum}
  If $P,Q$ are polytopes and $P\cup Q$ is a polytope, then
  \[
    (P\cup Q)+(P\cap Q)=P+Q.
  \]
\end{lemma}
\begin{proof}
  \uses{def:polytope, def:minkowski-sum}
  If $u\in P\cup Q$ and $v\in P\cap Q$, then $u+v\in P+Q$, choosing the
  two summands in the appropriate order.  This proves one inclusion.
  Conversely, take $p\in P$ and $q\in Q$.  The segment
  $z(t)=(1-t)p+tq$ lies in the convex set $P\cup Q$.  The closed subsets
  $\{t:z(t)\in P\}$ and $\{t:z(t)\in Q\}$ cover the connected interval
  $[0,1]$ and contain $0$ and $1$, respectively, so they intersect.  Choose
  $t$ in their intersection.  Then $z(t)\in P\cap Q$, while
  $z(1-t)\in P\cup Q$, and $z(t)+z(1-t)=p+q$.  This proves the reverse
  inclusion.
\end{proof}

\begin{lemma}[Lemma 9: full additivity]
  \label{lem:full-additivity}
  \uses{def:minkowski-sum, lem:two-cell-valuation}
  \notready
  Let $X,Q_1,\ldots,Q_m$ be polytopes with $X=\bigcup_{i=1}^mQ_i$.  For every
  nonempty $S\subseteq\{1,\ldots,m\}$ put $Q_S=\bigcap_{i\in S}Q_i$, omitting
  empty intersections, and let $\mathcal E$ and $\mathcal O$ be the nonempty
  even- and odd-cardinality index sets.  Then
  \[
    X+\sum_{S\in\mathcal E}Q_S=\sum_{S\in\mathcal O}Q_S.
  \]
\end{lemma}
% TODO: The cited full-additivity theorem is not proved.  The paper's route is
% to extend support functions to finite polyhedral unions using Euler
% characteristic, prove full additivity there, then recover the polytope
% identity.  Formalize those topological steps rather than naive induction.

\begin{theorem}[Finite triangulation of a polytope]
  \label{thm:polytope-triangulation}
  \uses{def:polytope}
  \notready
  Every polytope admits a finite polyhedral subdivision into simplices.
\end{theorem}
% TODO: Standard background invoked by the paper, but its proof is omitted and
% no ready general polytope-triangulation theorem was found locally.

\begin{lemma}[Support functions turn Minkowski sums into sums]
  \label{lem:support-minkowski-sum}
  \uses{def:support-function, def:minkowski-sum}
  For polytopes $P,Q$, $h_{P+Q}=h_P+h_Q$.
\end{lemma}
\begin{proof}
  \uses{def:support-function, def:minkowski-sum}
  For every $x$,
  \[
  h_{P+Q}(x)=\max_{p\in P,q\in Q}\langle x,p+q\rangle
   =\max_{p\in P}\langle x,p\rangle+
     \max_{q\in Q}\langle x,q\rangle.
  \]
\end{proof}

\begin{corollary}[Geometric signed-simplex representation]
  \label{cor:geometric-signed-simplex}
  \uses{thm:polytope-triangulation, lem:full-additivity,
    lem:support-minkowski-sum}
  The support function of every polytope is a finite signed sum of support
  functions of simplices and their nonempty common faces.
\end{corollary}
\begin{proof}
  \uses{thm:polytope-triangulation, lem:full-additivity,
    lem:support-minkowski-sum}
  Choose a finite triangulation.  Apply full additivity to its simplex cells,
  move the even-intersection sum to the other side in the Grothendieck group,
  and apply support functions.  Lemma~\ref{lem:support-minkowski-sum} changes
  every Minkowski sum into an ordinary sum, giving the asserted signed sum.
\end{proof}

\begin{lemma}[Lemma 10: cover criterion for maximum depth]
  \label{lem:cover-depth-bound}
  \uses{lem:full-additivity, lem:relu-depth-formal-difference,
    lem:max-newton-simplex}
  Suppose $\Delta_{n-1}=\bigcup_{i=1}^mQ_i$ and every nonempty intersection
  $Q_S=\bigcap_{i\in S}Q_i$ belongs to $\PolyC_k$.  Then $\MAX_n$ is
  computable with at most $k$ hidden layers.
\end{lemma}
\begin{proof}
  \uses{lem:full-additivity, lem:polytope-class-sum-closed,
    lem:relu-depth-formal-difference, lem:max-newton-simplex}
  Full additivity gives $\Delta_{n-1}+P=Q$, where $P$ is the finite sum of the
  even intersections and $Q$ the finite sum of the odd intersections.  Both
  lie in $\PolyC_k$ by closure under finite Minkowski sums.  Lemma 8 turns this
  formal difference into a $k$-hidden-layer representation of the support
  function of $\Delta_{n-1}$, which is $\MAX_n$.
\end{proof}

\begin{lemma}[The classes $\PolyC_k$ are closed under faces]
  \label{lem:polytope-class-face-closed}
  \uses{def:polytope-classes}
  \notready
  If $P\in\PolyC_k$ and $F$ is a nonempty face of $P$, then
  $F\in\PolyC_k$.
\end{lemma}
% TODO: Asserted without proof.  Formalize faces of Minkowski sums and convex
% joins, then induct on k.

\begin{definition}[Polyhedral subdivision]
  \label{def:polyhedral-subdivision}
  \uses{def:polytope}
  A finite family $(Q_i)_{i\in I}$ is a polyhedral subdivision of a polytope
  $X$ if $X=\bigcup_iQ_i$ and every nonempty intersection of cells is a face
  of each participating cell.
\end{definition}

\begin{corollary}[Subdivision criterion for maximum depth]
  \label{cor:subdivision-depth-bound}
  \uses{def:polyhedral-subdivision, lem:cover-depth-bound,
    lem:polytope-class-face-closed}
  If $\Delta_{n-1}$ has a polyhedral subdivision whose maximal cells belong to
  $\PolyC_k$, then $\MAX_n$ has hidden-layer depth at most $k$.
\end{corollary}
\begin{proof}
  \uses{lem:polytope-class-face-closed, lem:cover-depth-bound}
  Every nonempty multiple intersection is a face of a participating maximal
  cell, hence belongs to $\PolyC_k$ by face closure.  Lemma 10 applies.
\end{proof}

\section{Subdividing the four-simplex}

\begin{definition}[Construction: a coordinate tetrahedron]
  \label{constr:tetrahedron}
  \uses{def:polytope}
  In $\R^3$ put
  \[
  x_1=(-1,-1,-1),\quad x_2=(1,1,-1),\quad
  x_3=(-1,1,1),\quad x_4=(1,-1,1),
  \]
  and let $X=\operatorname{conv}\{x_1,x_2,x_3,x_4\}$.
\end{definition}

\begin{lemma}[The coordinate hull is a three-simplex]
  \label{lem:tetrahedron-simplex}
  \uses{constr:tetrahedron, def:standard-simplex}
  The points $x_1,x_2,x_3,x_4$ are affinely independent, so $X$ is affinely
  equivalent to $\Delta_3$.
\end{lemma}
\begin{proof}
  \uses{constr:tetrahedron}
  Relative to $x_1$, the other three difference vectors are
  $(2,2,0),(0,2,2),(2,0,2)$.  Their determinant is $16\neq0$, so they are
  linearly independent.  Mapping the four standard simplex vertices to the
  $x_i$ therefore extends to an affine isomorphism of their affine hulls and
  carries $\Delta_3$ to $X$.
\end{proof}

\begin{definition}[Edge midpoints]
  \label{def:tetrahedron-midpoints}
  \uses{constr:tetrahedron}
  For $1\leq i<j\leq4$, put $x_{ij}=(x_i+x_j)/2$.
\end{definition}

\begin{lemma}[Square projection]
  \label{lem:tetrahedron-projection-square}
  \uses{constr:tetrahedron}
  Orthogonal projection of $x_1,x_2,x_3,x_4$ onto $e_3^\perp$ gives the four
  vertices of the square $[-1,1]^2$.
\end{lemma}
\begin{proof}
  \uses{constr:tetrahedron}
  Projection deletes the third coordinate, giving
  $(-1,-1),(1,1),(-1,1),(1,-1)$, exactly the four square vertices.
\end{proof}

\begin{definition}[Construction: four rhombic-pyramid cells]
  \label{constr:rhombic-cells}
  \uses{def:tetrahedron-midpoints}
  Define four cells by their apex and four base vertices:
  \[
  \begin{array}{c|c}
  x_{12}&x_1,x_{14},x_{34},x_{13}\\
  x_{12}&x_2,x_{23},x_{34},x_{24}\\
  x_{34}&x_3,x_{13},x_{12},x_{23}\\
  x_{34}&x_4,x_{14},x_{12},x_{24}.
  \end{array}
  \]
  In each row $Q_i$ is the convex hull of the apex and the listed base, and
  $Z_i$ is the convex hull of the four base vertices.
\end{definition}

\begin{lemma}[The four bases are rhombi]
  \label{lem:rhombic-bases}
  \uses{constr:rhombic-cells}
  Each $Z_i$ is a parallelogram with four equal edge lengths, hence a rhombus.
\end{lemma}
\begin{proof}
  \uses{constr:tetrahedron, def:tetrahedron-midpoints}
  For the first base,
  $x_1+x_{34}=x_{14}+x_{13}$, so its diagonals bisect each other, while the
  adjacent difference vectors $(1,0,1)$ and $(0,1,1)$ both have length
  $\sqrt2$.  The coordinate symmetries swapping $1,2$, swapping $3,4$, and
  swapping the two pairs carry this computation to the other three bases.
\end{proof}

\begin{lemma}[A rhombus belongs to $\PolyC_1$]
  \label{lem:rhombus-class-one}
  \uses{def:polytope-classes, def:minkowski-sum}
  Every parallelogram, and hence every $Z_i$, belongs to $\PolyC_1$.
\end{lemma}
\begin{proof}
  \uses{def:polytope-classes}
  Translating if necessary, a parallelogram with edge vectors $u,v$ is
  $p+[0,u]+[0,v]$.  Each segment is the convex join of its two singleton
  endpoints, so the Minkowski sum of the two segments (with the translation
  absorbed into one segment's endpoints) is a finite sum of joins of
  $\PolyC_0$ members, hence is in $\PolyC_1$.
\end{proof}

\begin{lemma}[Each cell is a rhombic pyramid in $\PolyC_2$]
  \label{lem:rhombic-cells-class-two}
  \uses{constr:rhombic-cells, lem:rhombus-class-one,
    lem:polytope-class-monotone}
  For every $i$, $Q_i=p_i*Z_i$ for the listed apex $p_i$, and
  $Q_i\in\PolyC_2$.
\end{lemma}
\begin{proof}
  \uses{def:convex-join, def:polytope-classes,
    lem:rhombus-class-one, lem:polytope-class-monotone}
  The equality is the definition of the convex hull of the apex and base.
  The apex is in $\PolyC_0$ and therefore in $\PolyC_1$ by monotonicity, while
  $Z_i\in\PolyC_1$.  Their join is an allowed one-term expression in
  $\PolyC_2$.
\end{proof}

\begin{lemma}[The two coordinate cuts cover the tetrahedron]
  \label{lem:tetrahedron-cell-cover}
  \uses{constr:tetrahedron, constr:rhombic-cells,
    lem:tetrahedron-projection-square}
  \notready
  Cutting $X$ by the planes $e_1^\perp$ and $e_2^\perp$ yields exactly the
  four cells $Q_1,Q_2,Q_3,Q_4$, and $X=\bigcup_iQ_i$.
\end{lemma}
% TODO: The paper establishes this with Figures 1--2 and the cell table but
% omits a barycentric or half-space proof.  Verify the four sign regions.

\begin{lemma}[The tetrahedral cells meet face-to-face]
  \label{lem:tetrahedron-cell-faces}
  \uses{constr:rhombic-cells}
  \notready
  Every nonempty intersection of the cells $Q_i$ is a common face of the
  participating cells.
\end{lemma}
% TODO: Compute half-space descriptions and all nonempty incidences; the paper
% provides only the geometric figures.

\begin{corollary}[Rhombic-pyramid subdivision of $\Delta_3$]
  \label{cor:tetrahedron-subdivision}
  \uses{lem:tetrahedron-simplex, lem:rhombic-cells-class-two,
    lem:tetrahedron-cell-cover, lem:tetrahedron-cell-faces,
    def:polyhedral-subdivision}
  The four $Q_i$ give a polyhedral subdivision of a copy of $\Delta_3$ into
  rhombic pyramids belonging to $\PolyC_2$.
\end{corollary}
\begin{proof}
  \uses{lem:tetrahedron-simplex, lem:rhombic-cells-class-two,
    lem:tetrahedron-cell-cover, lem:tetrahedron-cell-faces}
  The affine-simplex identification, the cover, and the face-to-face property
  are exactly the subdivision axioms; class membership was proved separately.
\end{proof}

\begin{definition}[Construction: coning the subdivision to $\Delta_4$]
  \label{constr:delta-four-cells}
  \uses{cor:tetrahedron-subdivision, def:convex-join}
  Let $X$ be a tetrahedral facet of $\Delta_4$, let $x_5$ be its opposite
  vertex, and define $\widehat Q_i=x_5*Q_i$.
\end{definition}

\begin{lemma}[Reassociation of the coned cells]
  \label{lem:delta-four-cell-reassociate}
  \uses{constr:delta-four-cells, lem:convex-join-assoc}
  For $Q_i=p_i*Z_i$,
  $x_5*(p_i*Z_i)=(x_5*p_i)*Z_i$.
\end{lemma}
\begin{proof}
  \uses{lem:convex-join-assoc}
  This is associativity of convex join.
\end{proof}

\begin{lemma}[The coned cells belong to $\PolyC_2$]
  \label{lem:delta-four-cells-class-two}
  \uses{lem:delta-four-cell-reassociate, lem:rhombus-class-one,
    def:polytope-classes}
  Each $\widehat Q_i$ belongs to $\PolyC_2$.
\end{lemma}
\begin{proof}
  \uses{lem:delta-four-cell-reassociate, lem:rhombus-class-one,
    def:polytope-classes}
  The segment $x_5*p_i$ is a join of singleton polytopes and lies in
  $\PolyC_1$; the rhombus $Z_i$ also lies in $\PolyC_1$.  Their join is an
  allowed $\PolyC_2$ term.
\end{proof}

\begin{lemma}[The coned cells subdivide $\Delta_4$]
  \label{lem:delta-four-subdivision}
  \uses{constr:delta-four-cells, cor:tetrahedron-subdivision}
  The four cells $\widehat Q_i$ form a polyhedral subdivision of
  $\Delta_4=x_5*X$.
\end{lemma}
\begin{proof}
  \uses{cor:tetrahedron-subdivision, def:polyhedral-subdivision}
  Convex join with the common apex distributes over the union of the cells,
  so the coned cells cover $x_5*X$.  Because the apex lies outside the affine
  hull of $X$, intersections of cones are cones over the corresponding base
  intersections.  A cone over a common face is a common face of the coned
  cells.  Thus the face-to-face property is preserved.
\end{proof}

\begin{corollary}[Geometric two-layer realization of $\MAX_5$]
  \label{cor:max-five-geometric}
  \uses{lem:delta-four-cells-class-two, lem:delta-four-subdivision,
    cor:subdivision-depth-bound}
  The subdivision of $\Delta_4$ gives a two-hidden-layer realization of
  $\MAX_5$.
\end{corollary}
\begin{proof}
  \uses{lem:delta-four-cells-class-two, lem:delta-four-subdivision,
    cor:subdivision-depth-bound}
  Apply the subdivision criterion with $n=5$ and $k=2$.
\end{proof}

\begin{lemma}[Newton polytopes of the four positive full-dimensional terms]
  \label{lem:newton-p-cells}
  \uses{def:newton-polytope, def:p1, def:p2, def:p3, def:p4,
    constr:delta-four-cells}
  \notready
  Under the affine coordinate identification above, the unordered family of
  Newton polytopes of $P_1,P_2,P_3,P_4$ is the family of four full-dimensional
  cells $\widehat Q_1,\ldots,\widehat Q_4$.
\end{lemma}
% TODO: The paper asserts this correspondence without listing the affine map.
% Compare the convex hulls of the six coefficient vectors term by term.

\begin{lemma}[Newton polytopes of the correction terms]
  \label{lem:newton-correction-faces}
  \uses{def:newton-polytope, def:q, def:r13, def:r14, def:r23, def:r24,
    lem:newton-p-cells}
  \notready
  The Newton polytopes of $Q,R_{13},R_{14},R_{23},R_{24}$ are precisely the
  lower-dimensional faces that occur in the inclusion--exclusion identity for
  the four full-dimensional cells.
\end{lemma}
% TODO: The paper does not enumerate the face incidences.  Compute them from
% the coefficient-vector convex hulls before formalization.

\section{Higher-dimensional interpretation}

\begin{lemma}[A tetrahedron is a point--triangle join]
  \label{lem:tetrahedron-point-triangle}
  \uses{def:standard-simplex, def:convex-join}
  The simplex $\Delta_3$ is the convex join of any one vertex with the opposite
  triangular facet.
\end{lemma}
\begin{proof}
  The convex hull of all four vertices is, by associativity of convex hull,
  the convex hull of the chosen vertex together with the convex hull of the
  other three vertices.
\end{proof}

\begin{lemma}[Triangles have polytope complexity at most two]
  \label{lem:triangle-class-two}
  \uses{def:polytope-classes, lem:polytope-class-monotone}
  Every triangle belongs to $\PolyC_2$.
\end{lemma}
\begin{proof}
  \uses{def:polytope-classes, lem:polytope-class-monotone}
  A segment is the join of two singleton polytopes and belongs to
  $\PolyC_1$.  A third singleton also belongs to $\PolyC_1$ by monotonicity.
  Their join is the triangle and is an allowed term in $\PolyC_2$.
\end{proof}

\begin{proposition}[Three-to-one geometric replacement]
  \label{prop:geometric-three-to-one}
  \uses{lem:tetrahedron-point-triangle, lem:triangle-class-two,
    cor:tetrahedron-subdivision, lem:rhombus-class-one}
  The rhombic-pyramid subdivision replaces the triangular component of the
  point--triangle presentation of $\Delta_3$ (a join of three points in
  $\PolyC_2$) by a rhombus in $\PolyC_1$ inside each cell.
\end{proposition}
\begin{proof}
  \uses{lem:tetrahedron-point-triangle, lem:triangle-class-two,
    cor:tetrahedron-subdivision, lem:rhombus-class-one}
  The first two lemmas give the point--triangle presentation and its class.
  The explicit subdivision writes every cell as $p_i*Z_i$, and each $Z_i$ is
  a rhombus in $\PolyC_1$.  This is exactly the stated local replacement.
\end{proof}

\begin{proposition}[Omitted higher-dimensional geometric recursion]
  \label{remark:geometric-recursion-gap}
  \uses{prop:geometric-three-to-one, lem:claim-five}
  \notready
  A fully geometric iteration should formalize that this replacement reduces
  the number of joined objects by a factor of three per layer and recovers the
  $\log_3$ bound.  The paper deliberately omits the required theory of joins
  interacting with subdivisions; Claim~\ref{lem:claim-five} is its algebraic
  substitute.
\end{proposition}
% TODO: Supply the missing higher-dimensional join/subdivision theorem rather
% than promoting the heuristic sentence to a proof.

\chapter{Conclusions and external quantitative context}

\begin{theorem}[Universal approximation by one hidden layer]
  \label{thm:universal-approximation}
  \uses{def:relu-network}
  \notready
  On a bounded domain, every continuous function can be uniformly approximated
  arbitrarily well by one-hidden-layer networks under the hypotheses of the
  cited universal-approximation theorems; the nonpolynomial-activation version
  includes ReLU.
\end{theorem}
% TODO: The paper cites [Cyb89, HSW89, LLPS93]; state one precise theorem with
% its domain, norm, and activation hypotheses, then formalize or import it.

\begin{proposition}[Reported bound on the Wang--Sun term count]
  \label{prop:wang-sun-term-count}
  \uses{thm:wang-sun-representation}
  \notready
  If a CPWL function on $\R^n$ has $p$ distinct affine pieces represented
  locally, then the reported best asymptotic upper bound for the number $s$ of
  terms in Equation (1) is $s\in O(p^{n+1})$.
\end{proposition}
% TODO: The paper aggregates several references and gives no constants or
% exact quantifiers.  Pin down the precise source theorem before formalizing.

\begin{proposition}[Open depth questions]
  \label{remark:open-depth-questions}
  \uses{prop:max-five-depth, thm:main-cpwl}
  It remains open whether $\MAX_6$ has a two-hidden-layer realization and
  whether a constant number of hidden layers, possibly two, represents every
  CPWL function in every dimension.
\end{proposition}
\begin{proof}
  This records the paper's explicit open problems rather than a proposition
  requiring proof.
\end{proof}
